\documentclass[11pt,a4paper]{article}
\newcommand{\fichetype}{Corrigé détaillé}
% Thème A4 commun, librement modifiable.
\usepackage[T1]{fontenc}
\usepackage[utf8]{inputenc}
\usepackage[french]{babel}
\usepackage{lmodern,amsmath,amssymb,mathtools,tikz,pgfplots}
\usepackage[margin=2.2cm,headheight=15pt]{geometry}
\usepackage{xcolor,fancyhdr,enumitem,hyperref,needspace}
\usetikzlibrary{arrows.meta,positioning,calc}
\pgfplotsset{compat=1.18,width=11cm,height=6cm}
\definecolor{Ink}{HTML}{202124}
\definecolor{Blue}{HTML}{356B9E}
\definecolor{Muted}{HTML}{737579}
\color{Ink}
\renewcommand{\familydefault}{\sfdefault}
\setlength{\parindent}{0pt}
\setlength{\parskip}{6pt}
\setlist{itemsep=3pt,topsep=4pt}
\pagestyle{fancy}\fancyhf{}
\fancyhead[L]{\small\color{Muted}MATHS LYCÉE}
\fancyhead[R]{\small\color{Muted}2026--2027}
\fancyfoot[L]{\small\color{Muted}\fichetype}
\fancyfoot[R]{\small\thepage}
\renewcommand{\headrulewidth}{0.3pt}
\hypersetup{colorlinks=true,urlcolor=Blue,linkcolor=Blue}
\newcommand{\R}{\mathbb{R}}
\newcommand{\N}{\mathbb{N}}
\newcommand{\E}{\mathbb{E}}
\newcommand{\V}{\mathrm{V}}
\newcommand{\exercice}[3]{\Needspace{6\baselineskip}\vspace{0.5cm}{\color{Blue}\small\textbf{EXERCICE #1}\hfill #2}\par{\large\bfseries #3}\par}

\begin{document}
{\small\color{Blue}PREMIÈRE SPÉCIALITÉ}\par{\LARGE\bfseries Variations, extrema et optimisation}\par{\large Corrigé détaillé}\par\vspace{0.2cm}\hrule\vspace{0.4cm}
\par\noindent\begin{minipage}{\linewidth}\exercice{1}{Application}{Une parabole}
{\small\color{Muted}\textbf{Énoncé.} Étudier $f(x)=x^2-6x+5$ sur $[0,5]$ et donner ses extrema.}
\par\medskip
\textbf{Résolution.}\par
$f\prime=2x-6$ est négative avant $3$, positive après. Ainsi $f$ décroît de $f(0)=5$ à $f(3)=-4$, puis croît jusqu’à $f(5)=0$. Minimum $-4$ en $3$, maximum $5$ en $0$.
\end{minipage}\par
\par\noindent\begin{minipage}{\linewidth}\exercice{2}{Application}{Un cube}
{\small\color{Muted}\textbf{Énoncé.} Étudier les variations de $x^3+2x$ sur $\R$.}
\par\medskip
\textbf{Résolution.}\par
La dérivée $3x^2+2$ est strictement positive pour tout réel. La fonction est strictement croissante sur $\R$. Il n’y a pas d’extremum local.
\end{minipage}\par
\par\noindent\begin{minipage}{\linewidth}\exercice{3}{Entraînement}{Deux intervalles}
{\small\color{Muted}\textbf{Énoncé.} Étudier $1/x$ sur son domaine. Peut-on dire qu’elle est décroissante sur tout son domaine ?}
\par\medskip
\textbf{Résolution.}\par
La dérivée est $-1/x^2<0$. La fonction décroît sur chacun des intervalles $]-\infty,0[$ et $]0,+\infty[$. Elle ne décroît pas sur leur réunion : $-1<1$ mais $f(-1)=-1<f(1)=1$. Le théorème exige un intervalle.
\end{minipage}\par
\par\noindent\begin{minipage}{\linewidth}\exercice{4}{Entraînement}{Rectangle}
{\small\color{Muted}\textbf{Énoncé.} Un rectangle a un périmètre de $20$ m. Quelle aire maximale peut-il avoir ?}
\par\medskip
\textbf{Résolution.}\par
Pour une largeur $x\in]0,10[$, la longueur est $10-x$. L’aire vaut $A=10x-x^2$, donc $A\prime=10-2x$. Elle est maximale en $x=5$, avec $A(5)=25$ m$^2$. Le rectangle optimal est un carré.
\end{minipage}\par
\par\noindent\begin{minipage}{\linewidth}\exercice{5}{Application}{Parité}
{\small\color{Muted}\textbf{Énoncé.} Étudier la parité de $x^4-3x^2$, $x^3+x$ et $x^2+x+1$.}
\par\medskip
\textbf{Résolution.}\par
Les domaines sont symétriques. Remplacer $x$ par $-x$ laisse la première expression inchangée : elle est paire. La deuxième devient son opposée : elle est impaire. Pour la troisième, $f(1)=3$, $f(-1)=1$ : elle n’est ni paire ni impaire.
\end{minipage}\par
\par\noindent\begin{minipage}{\linewidth}\exercice{6}{Entraînement}{Une inégalité}
{\small\color{Muted}\textbf{Énoncé.} Montrer que $x+1/x\ge2$ pour $x>0$ par une étude de variations.}
\par\medskip
\textbf{Résolution.}\par
Poser $f=x+1/x$. Sa dérivée est $(x^2-1)/x^2$. Le dénominateur est positif et le numérateur change de signe en $1$. Le minimum vaut $f(1)=2$. L’égalité a lieu seulement en $1$.
\end{minipage}\par
\par\noindent\begin{minipage}{\linewidth}\exercice{7}{Approfondissement}{Production entière}
{\small\color{Muted}\textbf{Énoncé.} Un bénéfice fictif est $B(n)=-2n^2+14n-5$ pour un entier $0\le n\le8$. Chercher le maximum.}
\par\medskip
\textbf{Résolution.}\par
La fonction continue associée a une dérivée $-4x+14$, nulle en $3,5$. Elle croît puis décroît. On compare $B(3)=19$ et $B(4)=19$, qui sont les deux maxima entiers. Arrondir sans comparer aurait fait perdre une solution.
\end{minipage}\par
\par\noindent\begin{minipage}{\linewidth}\exercice{8}{Approfondissement}{Contre-exemple}
{\small\color{Muted}\textbf{Énoncé.} Expliquer pourquoi $f\prime(a)=0$ n’implique pas un extremum. Donner un exemple où un maximum existe au bord avec une dérivée non nulle.}
\par\medskip
\textbf{Résolution.}\par
Pour $x^3$ en zéro, la dérivée est nulle mais la fonction croît avant et après. Pour $f(x)=x$ sur $[0,1]$, le maximum vaut $1$ au bord $x=1$, alors que la dérivée vaut $1$. La condition de dérivée nulle concerne les extrema intérieurs d’une fonction dérivable.
\end{minipage}\par
\end{document}
